%=======================================================================
%  CHAPTER 6 - IIR FILTER DESIGN  (6 hrs)
%  Every order, cut-off, pole, H(s) and H(z) quoted here is produced by
%  iir.py, which self-tests against the published Oppenheim examples 7.2,
%  7.3 and 7.6 before it prints, and is asserted again by verify.py ch6.
%  Sources: Chapterwise Ch6 "Design of IIR Filters", 36 pp, the closest
%           thing the subject has to a chapter-native text: pp.2-15 impulse
%           invariance with the Butterworth and Chebyshev worked examples,
%           pp.16-22 bilinear, pp.23-30 frequency transformation of low-pass
%           IIR filters, pp.30-36 structures (which is chapter 4 material).
%           SPP Sir deck 6 "IIR _ FIR" pp.1-105 (pp.106-155 are chapter 5):
%           pp.19-25 the three reasons for designing through an analog
%           prototype, pp.36-47 impulse invariance, pp.48-56 bilinear and
%           the comparison, pp.57-62 the prototype families, pp.100-105
%           the frequency-transformation table.
%           DSAP All Chapter Summary pp.72-85 for the LOCAL convention: the
%           order and cut-off formulas written in dB attenuation, which is
%           how every one of the 45 papers states its specification, and the
%           digital spectral-transformation table with its alpha and k.
%  NOTE the syllabus file's chapter 6 is corrupt (it repeats the
%  RF-Microwave antenna topics). The reconstruction used here is in
%  DSAP/NOTES-BUILD-PROGRESS.md 5.1 and matches what the papers ask.
%  Numericals live in ch6-num.tex.
%=======================================================================
\section{IIR Filter Design}

{\color{sub}\footnotesize 6 hrs \gap$\cdot$\gap \mk{50 questions across the 45 papers}
\gap$\cdot$\gap worth about \mk{13 of 80 marks} \gap$\cdot$\gap 4 syllabus sub-topics
\gap$\cdot$\gap \mk{2nd of the 8 chapters} by weight, behind only the DFT.}

\lead{Read this first:}
\begin{itemize}
  \item \mk{36 of the 50 questions are one calculation}: design an analog Butterworth
        low-pass filter, then map it to digital. Learn that one procedure and the
        chapter is mostly done.
  \item The whole subject of IIR design is \textbf{two sentences}. There is no good
        direct method in the $z$ domain, but analog filter design is a solved,
        closed-form art. So \textbf{design in $s$, then map $s\to z$}.
  \item \textbf{Only two maps exist} in this syllabus, and they split the chapter:
  \begin{itemize}
    \item \textbf{Bilinear} \tS{36}: $s=\dfrac{2}{T}\dfrac{1-z^{-1}}{1+z^{-1}}$. One to
          one, so \textbf{no aliasing}, but the frequency axis is \textbf{warped}, so
          the edges must be pre-warped first.
    \item \textbf{Impulse invariance} \tS{11}: $h[n]=h_a(nT)$. The frequency axis is
          \textbf{linear}, so no warping, but sampling an unbandlimited analog response
          \textbf{aliases}.
  \end{itemize}
  \item \mk{Everything else is the analog half}: which prototype (Butterworth or
        Chebyshev), what order, what cut-off. That part is identical for both maps.
  \item \mk{The sampling period cancels.} When the spec is given in rad/sample, pre-warp
        by $\Omega=\frac{2}{T}\tan\frac{\omega}{2}$ and map back by
        $s=\frac{2}{T}\frac{1-z^{-1}}{1+z^{-1}}$, and
        \textbf{every $2/T$ divides out}: $H(z)$ is the same for $T=1$, $T=2$ or
        $T=0.1$. Nine papers give a sampling frequency purely as a distraction.
        \mk{Set $T=1$ and write one line saying why} (\S6.5).
  \item The other 14 questions are bookwork: why transform at all, the two methods
        compared, pre-warping, and spectral transformation. All short, all below.
\end{itemize}

\hr

\T{6.1 \quad Why design through an analog prototype}

\Q{Why do we use transformation of continuous time filter design methods to design
discrete time IIR filters?}{\m{3+6}\m{3}}
{\color{sub}\footnotesize \tP{2} \yr{\texttt{70 Ma} \m{3+6}, 73 Shr \m{3}}
\gap$\cdot$\gap \mk{three reasons, one line each, then move on}}

\sbs{%
\lead{The three reasons}
\begin{enumerate}
  \item \textbf{The analog art is mature.} Continuous-time IIR design is a century old
        and its results are tabulated. Reusing it is free.
  \item \textbf{Closed-form formulas exist.} Butterworth, Chebyshev and elliptic
        prototypes all have closed-form order, cut-off and pole expressions, so the
        design is arithmetic, not search.
  \item \textbf{No such formulas exist in the $z$ domain.} The same approximation
        methods applied directly to a discrete-time filter do \emph{not} give closed
        forms, because a digital frequency response is \mk{periodic} and an analog one
        is not.
\end{enumerate}}{%
\lead{The three steps that follow from it}
\begin{enumerate}
  \item \textbf{Translate the spec} from digital $\omega$ to analog $\Omega$, by the
        rule that belongs to the map being used.
  \item \textbf{Design $H_a(s)$}: choose the prototype family, find $N$ and
        $\Omega_c$, place the poles, write $H_a(s)$.
  \item \textbf{Map back} to $H(z)$ with the same rule.
\end{enumerate}
\lead{What the map must preserve}
\begin{itemize}
  \item The \textbf{$j\Omega$ axis} of the $s$ plane must land on the \textbf{unit
        circle} of the $z$ plane, or the frequency response is not preserved.
  \item The \textbf{left half plane} must land \textbf{inside} the unit circle, or a
        stable analog filter becomes an unstable digital one.
  \item \mk{Both maps below satisfy both conditions.} That is the whole reason there
        are only two of them.
\end{itemize}}

\hr

\T{6.2 \quad The analog prototypes}

\Q{Compare the Butterworth approximation with the Chebyshev approximation}{\m{1+10}\m{3}}
{\color{sub}\footnotesize \tP{2} \yr{\textbf{\texttt{76 Bh}} \m{1+10}, \textbf{82 Bh}}
\gap$\cdot$\gap \mk{a table, not prose}}

\sbsr{0.5}{%
\lead{Butterworth: maximally flat}
\[ |H_a(j\Omega)|^2=\frac{1}{1+\left(\dfrac{\Omega}{\Omega_c}\right)^{2N}} \]
\begin{itemize}
  \item \textbf{All pole.} Every zero is at $\Omega\to\infty$.
  \item $|H_a(j0)|=1$ for every $N$, and
        $|H_a(j\Omega_c)|=\frac{1}{\sqrt2}$, i.e. \mk{$\Omega_c$ is the 3 dB point
        whatever the order}.
  \item \textbf{Monotonic} in both bands: no ripple anywhere.
  \item Poles lie on a \textbf{circle} of radius $\Omega_c$, equally spaced, none on
        the $j\Omega$ axis.
\end{itemize}}{%
\lead{Chebyshev type I: equiripple passband}
\[ |H_a(j\Omega)|^2=\frac{1}{1+\varepsilon^2V_N^2(\Omega/\Omega_c)} \]
\begin{itemize}
  \item $V_N(x)=\cos(N\cos^{-1}x)$ for $|x|\le1$, which oscillates between $\pm1$, and
        $\cosh(N\cosh^{-1}x)$ for $x>1$, which climbs fast.
  \item So the passband \textbf{ripples} between $1$ and $1/\sqrt{1+\varepsilon^2}$, and
        the stopband falls monotonically and \textbf{steeply}.
  \item \textbf{Type II} puts the ripple in the stopband instead, by using
        $V_N^{-1}$; \textbf{elliptic} ripples in both and is steeper still.
  \item \mk{$|H_a(j0)|=1$ for $N$ odd but $1/\sqrt{1+\varepsilon^2}$ for $N$ even.}
        An even-order Chebyshev filter does \textbf{not} have unity DC gain.
\end{itemize}}

\lead{The comparison, as the answer}
\par\smallskip
\begin{tabular}{@{}L{2.6cm}L{5.4cm}L{6.9cm}@{}}
\toprule
\hrow \thead{} & \thead{Butterworth} & \thead{Chebyshev I} \\
Passband & monotonic, maximally flat & equiripple between $1$ and
  $1/\sqrt{1+\varepsilon^2}$ \\
Stopband & monotonic & monotonic \\
Transition & \textbf{widest} of the three families & narrower for the same order \\
Order for a given spec & \textbf{highest} & \mk{lower, typically by a third} \\
Poles lie on & a circle, radius $\Omega_c$ & an \textbf{ellipse}, semi-axes
  $a\Omega_c$ and $b\Omega_c$ \\
Phase & most nearly linear of the three & worse, and worst near the band edge \\
Use when & flatness matters more than sharpness & a short filter matters more than
  flatness \\
\bottomrule
\end{tabular}
\par\smallskip
{\color{sub}\footnotesize The number to quote comes from the 45 papers' own recurring
spec ($\omega_p=0.2\pi$, $\omega_s=0.3\pi$, 1 dB, 15 dB), asked by 73 Shr, 77 Ch and
76 Bh: \mk{Butterworth needs $N=6$, Chebyshev only $N=4$}. That single sentence answers
"compare" better than a paragraph. \added\ The local notes offer a second example,
$\Omega_p=1$, $\Omega_s=3.33$, $\alpha_p=0.3$ dB, $\alpha_s=22$ dB, and print
"$n=5$ for butterworth, $n=3$ for chebyshev". \textbf{The Butterworth figure is wrong}:
the ratio is $157.489/0.071519=2202.05$, its log is $3.3429$, and
$2\log(3.33)=1.0449$, so $N=3.1992$ and the answer is \mk{4, not 5}. The Chebyshev 3 is
right. Use the paper-derived example above and this one only with the correction.}

\Q{The design formulas, in the form every paper's specification uses}{\m{core}}
{\color{sub}\footnotesize \mk{these four lines carry 36 questions}
\gap$\cdot$\gap memorise them exactly}

\sbsr{0.44}{%
\figT{c6_tolerance.png}
{\color{sub}\footnotesize The tolerance scheme that every paper states in words.
$\delta_1$ is the passband deviation, $\delta_2$ the stopband gain, and \mk{only the
magnitude is ever specified}: an IIR filter's phase is whatever the poles give.}}{%
\lead{Attenuation is always positive dB}
\[ \alpha=-20\log_{10}|H| \]
\begin{itemize}
  \item "maximum deviation $1$ dB below 0 dB gain" $\Rightarrow\alpha_p=1$.
  \item "maximum gain $-15$ dB in the stopband" $\Rightarrow\alpha_s=15$.
  \item a gain $0.89125$ given directly $\Rightarrow\alpha_p=-20\log_{10}0.89125=1$ dB.
  \item a ripple $\delta_p=\delta_1=0.11$ $\Rightarrow$ the edge gain is
        $1-\delta_p=0.89$, so $\alpha_p=1.0122$ dB; but $\delta_s=\delta_2=0.21$ is a
        gain already, so $\alpha_s=13.5556$ dB.
  \item \mk{Those four phrasings are the same spec four ways.} Converting first, before
        anything else, is what makes the rest mechanical.
\end{itemize}}

\sbs{%
\lead{Order and cut-off}
\[ N\ \ge\ \frac{\log_{10}\!\left[\dfrac{10^{0.1\alpha_s}-1}
                                        {10^{0.1\alpha_p}-1}\right]}
                {2\log_{10}(\Omega_s/\Omega_p)}
   \qquad\text{(Butterworth)} \]
\[ N\ \ge\ \frac{\cosh^{-1}\!\sqrt{\dfrac{10^{0.1\alpha_s}-1}
                                         {10^{0.1\alpha_p}-1}}}
                {\cosh^{-1}(\Omega_s/\Omega_p)}
   \qquad\text{(Chebyshev)} \]
\[ \Omega_c=\frac{\Omega_p}{\left(10^{0.1\alpha_p}-1\right)^{1/2N}},\qquad
   \varepsilon=\sqrt{10^{0.1\alpha_p}-1} \]
\begin{itemize}
  \item \mk{Always round $N$ UP.} The fractional $N$ meets the spec exactly and is not
        realisable.
  \item That $\Omega_c$ makes the \textbf{passband} exact and leaves the stopband
        over-satisfied. Using $\Omega_s$ and $\alpha_s$ instead makes the stopband
        exact. Both are legal; \textbf{state which you used}.
  \item Chebyshev needs no separate $\Omega_c$: the passband is equiripple right up to
        $\Omega_p$, so \mk{$\Omega_c=\Omega_p$}.
\end{itemize}}{%
\lead{Butterworth poles}
\[ s_k=\Omega_c\,e^{\,j\pi\frac{N+2k+1}{2N}},\qquad k=0,1,\dots,N-1 \]
\begin{itemize}
  \item These are the $N$ left-half-plane roots of
        $1+(s/j\Omega_c)^{2N}=0$; the other $N$ belong to $H_a(-s)$ and are discarded.
  \item Equally spaced by $\pi/N$ on the circle of radius $\Omega_c$, symmetric about
        the real axis, \mk{never on the $j\Omega$ axis}.
  \item A real pole at $-\Omega_c$ appears \textbf{only for odd $N$}.
  \item[] \[ H_a(s)=\frac{\Omega_c^{\,N}}{\prod_{k}(s-s_k)} \]
  \item The numerator $\Omega_c^N$ is fixed by $H_a(0)=1$, not chosen.
\end{itemize}}

\sbs{%
\lead{Chebyshev poles}
\[ \alpha=\frac1\varepsilon+\sqrt{1+\frac1{\varepsilon^2}},\quad
   a=\frac{\alpha^{1/N}-\alpha^{-1/N}}{2},\quad
   b=\frac{\alpha^{1/N}+\alpha^{-1/N}}{2} \]
\[ s_k=-a\Omega_c\sin\theta_k+j\,b\Omega_c\cos\theta_k,\quad
   \theta_k=\frac{(2k+1)\pi}{2N} \]
\begin{itemize}
  \item The poles sit on the \textbf{ellipse} $\dfrac{\sigma^2}{a^2\Omega_c^2}
        +\dfrac{\Omega^2}{b^2\Omega_c^2}=1$, at the same angles a Butterworth filter
        would use, squashed towards the $j\Omega$ axis.
  \item $b>1>a$ always, so the ellipse is \textbf{taller than wide}: the poles crowd
        the imaginary axis, which is what sharpens the cut-off.
\end{itemize}}{%
\lead{Chebyshev gain, and the even-order trap}
\begin{itemize}
  \item $H_a(s)=\dfrac{K}{\prod_k(s-s_k)}$ with $K=\prod_k(-s_k)$ for \textbf{odd} $N$,
        which makes $H_a(0)=1$.
  \item For \textbf{even} $N$ the equiripple passband starts at its \emph{lower} edge, so
        $H_a(0)=\dfrac{1}{\sqrt{1+\varepsilon^2}}$ and \mk{$K$ is divided by
        $\sqrt{1+\varepsilon^2}$}.
  \item Worked instance: 72 Ma has $\varepsilon=1$, so its $N=2$ filter has a DC gain of
        $1/\sqrt2=0.707$, \mk{not 1}. It looks like an arithmetic error and is not.
  \item $\varepsilon$ and $\Omega_c$ come entirely from the passband. \textbf{Only $N$}
        uses the stopband at all.
\end{itemize}}

\hr

\T{6.3 \quad The impulse invariance method}

\Q{Describe digital Butterworth filter design using the impulse invariance technique.
What are the limitations of the impulse invariance technique?}{\m{15}\m{12+3}}
{\color{sub}\footnotesize \tS{11} \yr{\textbf{69 Bh} \m{15}, \textbf{78 Bh} \m{12+3},
\textbf{72 Ch} \m{3+12}, \textbf{\texttt{81 Ch}}, \textbf{\texttt{80 Ch}},
\textbf{\texttt{69 Bh}} \m{12}, 70 Asa \m{13}, \textbf{74 Ch} \m{12},
\textbf{67 Mng} \m{15}} \gap$\cdot$\gap \mk{the limitations are half the marks}}

{\color{sub}\footnotesize In one line: \mk{make the digital impulse response the sampled
analog one}, $h[n]=h_a(nT)$. Everything below follows from that single sentence, and so
does everything wrong with the method.}

\sbsr{0.54}{%
\lead{The idea, and the whole derivation}
\begin{itemize}
  \item \textbf{Sample the analog impulse response}: $h[n]=h_a(nT)$. Nothing else is
        assumed.
  \item Expand $H_a(s)$ in partial fractions over its $N$ simple poles $p_k$:
  \item[] \[ H_a(s)=\sum_{k=1}^{N}\frac{A_k}{s-p_k}
             \ \Longrightarrow\ h_a(t)=\sum_k A_k e^{p_kt}u(t) \]
  \item Sample it and take the $z$ transform, which is just a geometric series in
        $e^{p_kT}z^{-1}$:
  \item[] \[ H(z)=\sum_{n=0}^{\infty}\!\Big(\sum_k A_ke^{p_knT}\Big)z^{-n}
                 =\boxed{\sum_{k=1}^{N}\frac{A_k}{1-e^{p_kT}z^{-1}}} \]
  \item So the rule is \mk{term by term}:
        $\dfrac{1}{s-p_k}\longrightarrow\dfrac{1}{1-e^{p_kT}z^{-1}}$. That is the only
        thing to remember.
  \item \textbf{Spec translation is linear}: $\Omega=\omega/T$. \mk{No pre-warping.}
  \item \textbf{The $T$ convention.} Some texts set $h[n]=T\,h_a(nT)$ instead, which
        multiplies every $A_k$ by $T$. Without it the filter's gain comes out near
        $1/T$, so with $T\ne1$ you must say which you used. \mk{With $T=1$ s the two
        are identical}, and 8 of these 11 papers set $T=1$.
\end{itemize}}{%
\lead{What the map does to the planes}
\begin{itemize}
  \item Each pole moves as $\;p_k\to e^{p_kT}$, so
        $|z|=e^{\sigma_kT}$ and $\angle z=\Omega_kT$.
  \item $\sigma_k<0\Rightarrow|z|<1$: \mk{a stable analog filter always gives a stable
        digital one}.
  \item But $z=e^{sT}$ is \textbf{many to one}: every horizontal strip of width
        $2\pi/T$ in the $s$ plane maps onto the whole $z$ plane.
  \item \mk{This applies to the POLES only, not to $H$.} $H(z)\ne H_a(s)|_{s}$ under
        any substitution; only the pole positions transform that simply.
\end{itemize}}

\sbsr{0.54}{%
\lead{The limitations (say all four)}
\begin{enumerate}
  \item \textbf{Aliasing.} Sampling gives
        $H(e^{j\omega})=\frac1T\sum_{k}H_a\!\left(j\frac{\omega+2\pi k}{T}\right)$.
        No practical analog prototype is band-limited, so the tails fold back and the
        response is \mk{never exactly the analog one}.
  \item \textbf{Low pass and band pass only.} A high-pass or band-stop prototype has
        an infinite high-frequency tail, so it aliases catastrophically.
        \mk{Impulse invariance cannot make a high-pass filter.}
  \item \textbf{Needs a partial fraction expansion} of $H_a(s)$, and it assumes the
        poles are \textbf{simple}. Repeated poles need a separate rule.
  \item \textbf{Gain depends on $T$}, so the result must be rescaled.
\end{enumerate}}{%
\lead{The spec translation, which is all of step 1}
\begin{itemize}
  \item Take the tolerance scheme of \S6.2 and \mk{divide every frequency edge by $T$}.
        Nothing else changes: the bands $1-\delta_1$ and $\delta_2$ are untouched.
  \item The prototype is then specified only up to $\pi/T$, above which it is
        \emph{assumed} negligible. \mk{That assumption is exactly what aliasing
        violates}, which is why the design is checked at the end rather than trusted.
  \item Contrast \S6.4, where the axis is \textbf{bent} rather than merely rescaled.
  \item Practical consequence: \textbf{over-design the stopband}. Extra attenuation
        there is what buys back the margin that the folded tails eat.
\end{itemize}}

\hr

\T{6.4 \quad The bilinear transformation}

\Q{Explain briefly the bilinear transformation method of IIR filter design, with its
advantages and disadvantages}{\m{3}\m{3+6}\m{2+10}}
{\color{sub}\footnotesize \tS{36} \yr{73 Shr \m{3}, \texttt{70 Ma} \m{3+6},
81 Ba \m{2+10}, and the 32 design papers in \S6.5}}

\sbsr{0.52}{%
\lead{Where the formula comes from}
\begin{itemize}
  \item Take the simplest analog filter, $H_a(s)=\dfrac{b}{s+a}$, i.e.
        $\dfrac{dy}{dt}+ay(t)=bx(t)$.
  \item Integrate over one sample interval and apply the \textbf{trapezoidal rule}:
  \item[] \[ y(nT)-y(nT-T)=\frac{T}{2}\big[g(nT)+g(nT-T)\big],\quad g=bx-ay \]
  \item Take $z$ transforms and solve. The result is $H_a(s)$ with
  \item[] \[ \boxed{\;s=\frac{2}{T}\,\frac{1-z^{-1}}{1+z^{-1}}\;} \]
  \item So \mk{the bilinear transformation IS trapezoidal integration}, which is why it
        is stable and one to one. That sentence is worth stating: it is the only
        derivation in the chapter and papers ask for it.
  \item Inverting, $\;z=\dfrac{1+sT/2}{1-sT/2}$.
\end{itemize}}{%
\lead{What it does to the planes}
\begin{itemize}
  \item Put $s=\sigma+j\Omega$. Then
  \item[] \[ |z|^2=\frac{(1+\sigma T/2)^2+(\Omega T/2)^2}
                        {(1-\sigma T/2)^2+(\Omega T/2)^2} \]
  \item $\sigma<0\Rightarrow|z|<1$, $\sigma=0\Rightarrow|z|=1$,
        $\sigma>0\Rightarrow|z|>1$. \mk{The entire left half plane maps into the unit
        disc, exactly once.}
  \item So stability is preserved, and \textbf{there is no aliasing}: the map is one to
        one, not many to one.
  \item On the axis, $s=j\Omega$ gives $z=e^{j\omega}$ with
  \item[] \[ \omega=2\tan^{-1}\!\frac{\Omega T}{2},\qquad
             \boxed{\;\Omega=\frac{2}{T}\tan\frac{\omega}{2}\;} \]
\end{itemize}}

\Q{What is frequency warping, what is pre-warping, and why is it necessary}{\m{3+12}\m{3}}
{\color{sub}\footnotesize \tF{4} \yr{\textbf{78 Bh} \m{3}, \textbf{82 Bh} \m{3+12},
\textbf{72 Ch} \m{3}, 79 Ba \m{3}} \gap$\cdot$\gap \mk{the figure below is the answer}}

\sbsr{0.45}{%
\figT{c6_warping.png}
{\color{sub}\footnotesize Digital $\omega$ against analog $\Omega T_d$. The curve is
$\omega=2\tan^{-1}(\Omega T_d/2)$: a straight line through the origin at low frequency,
flattening to the horizontal asymptotes $\pm\pi$.}}{%
\lead{Warping}
\begin{itemize}
  \item $\Omega$ runs over all of $(-\infty,\infty)$ but $\omega$ only over
        $(-\pi,\pi)$, so \mk{the infinite analog axis is compressed into a finite
        digital one}. That compression is the warping.
  \item It is \textbf{almost linear for small $\omega$} ($\tan x\approx x$) and
        \textbf{severely non-linear near $\omega=\pi$}.
  \item Consequence: a shape defined by more than its band edges, such as a
        differentiator or an arbitrary magnitude curve, comes out \mk{distorted}.
        Bilinear design is only safe for \textbf{piecewise-constant} magnitude
        responses: low pass, high pass, band pass, band stop.
\end{itemize}
\lead{Pre-warping}
\begin{itemize}
  \item Because the map bends the axis, designing the analog filter at the \emph{stated}
        edges puts the digital edges somewhere else.
  \item So \textbf{push the edges through the same bend first}:
        $\Omega_p=\frac{2}{T}\tan\frac{\omega_p}{2}$,
        $\Omega_s=\frac{2}{T}\tan\frac{\omega_s}{2}$.
  \item The warp then undoes the pre-warp exactly, and \mk{the critical frequencies land
        where the specification asked}. Everything between them is still bent, which is
        why only the edges can be guaranteed.
  \item \mk{Skipping the pre-warp is the single most common lost mark in the chapter.}
\end{itemize}}

\hr

\T{6.5 \quad Impulse invariance against bilinear}

\Q{Compare the impulse invariance method with the bilinear transformation method}%
{\m{12+3}\m{11+3}\m{3+12}\m{2+10}}
{\color{sub}\footnotesize \tF{6} \yr{\textbf{81 Bh} \m{12+3}, 79 Ba \m{11+3},
82 Ba \m{12+3}, \textbf{72 Ch} \m{3+12}, 81 Ba \m{2+10}, \texttt{70 Ma} \m{3+6}}
\gap$\cdot$\gap \mk{asked as a rider on a design question six times}}

\par\smallskip
\begin{tabular}{@{}L{3.0cm}L{5.7cm}L{6.2cm}@{}}
\toprule
\hrow \thead{} & \thead{Impulse invariance} & \thead{Bilinear} \\
Basis & $h[n]=h_a(nT)$, a time-domain match &
  trapezoidal integration, a frequency-domain map \\
$s\to z$ & poles only: $p_k\to e^{p_kT}$ &
  $s=\dfrac{2}{T}\dfrac{1-z^{-1}}{1+z^{-1}}$ throughout \\
Frequency map & $\Omega=\omega/T$, \textbf{linear} &
  $\Omega=\frac{2}{T}\tan\frac{\omega}{2}$, \textbf{non-linear} \\
Mapping type & many to one & \mk{one to one} \\
Aliasing & \mk{yes}, and unavoidable & \mk{none} \\
Warping & none, shape preserved & yes, needs pre-warping \\
Filter types & low pass, band pass only & \textbf{all four} \\
Non-constant shapes & can do a band-limited differentiator & cannot \\
Needs & partial fractions, simple poles & algebraic substitution only \\
Order needed & higher, to absorb the aliasing & lower \\
\bottomrule
\end{tabular}

\sbs{%
\lead{The one-sentence verdict}
\begin{itemize}
  \item \mk{Bilinear trades shape fidelity for the absence of aliasing; impulse
        invariance does the reverse.}
  \item Since almost every filter wanted is piecewise constant, where shape fidelity
        between the edges is worth nothing, \textbf{bilinear wins almost always}, and
        36 of this chapter's 50 questions use it.
  \item The exception is a filter whose \emph{shape} matters across the band. Only
        impulse invariance can make a band-limited differentiator.
\end{itemize}}{%
\lead{Why $T$ cancels in a bilinear design \added}
\begin{itemize}
  \item Pre-warping gives $\Omega_p=\frac{2}{T}\tan\frac{\omega_p}{2}$, so
        $\Omega_s/\Omega_p$ is \textbf{free of $T$} and so is $N$.
  \item $\Omega_c$ then carries one factor $\frac{2}{T}$, so every pole does, so
        $H_a(s)$ is a function of $sT/2$ alone.
  \item Substituting $s=\frac{2}{T}\frac{1-z^{-1}}{1+z^{-1}}$ cancels it.
        \mk{$H(z)$ is identical for every $T$.}
  \item So a paper's "sampling frequency 0.5 Hz" changes nothing in the answer. Use
        $T=1$, and \textbf{say in one line that it cancels} to show it was not ignored.
  \item \mk{This is false for impulse invariance}, where $\Omega=\omega/T$ scales the
        poles but $e^{p_kT}$ does not cancel it.
\end{itemize}}

\hr

\T{6.6 \quad Spectral transformation in the digital domain}

\Q{Why is spectral transformation required, and what are its features and parameters
for low pass to high pass conversion}{\m{2}\m{4}\m{11+4}}
{\color{sub}\footnotesize \tF{6} \yr{70 Asa \m{2}, 80 Ba \m{4}, \textbf{70 Ch} \m{5},
\textbf{79 Bh}, \textbf{71 Ch}, \textbf{69 Ch} \m{11+4}}}

\sbs{%
\lead{Why it is required}
\begin{itemize}
  \item Every design procedure above produces a \textbf{low pass} filter. High pass,
        band pass and band stop have to come from somewhere.
  \item The analog route (transform $H_a(s)$ into a high pass, \emph{then} discretise)
        works for bilinear but \mk{fails for impulse invariance}, because a high-pass
        analog filter is not band-limited and aliases.
  \item So do it the other way: \textbf{discretise the low pass first}, then transform
        \textbf{inside the $z$ domain}. Aliasing is then impossible, because nothing is
        sampled.
  \item \mk{For bilinear designs the two routes give the identical filter}; for impulse
        invariance only the digital route exists. That is the answer to "why required".
\end{itemize}}{%
\lead{What the map must satisfy}
\begin{itemize}
  \item $p^{-1}=G(z^{-1})$ must be \textbf{rational} in $z^{-1}$.
  \item It must take the unit circle to the unit circle, and the inside to the inside,
        so that \textbf{stability survives}.
  \item The general form that does both is an \textbf{all-pass} function,
        $\prod_k\dfrac{z^{-1}-\alpha_k}{1-\alpha_kz^{-1}}$ with every $|\alpha_k|<1$.
  \item \mk{$N$ all-pass factors give an $N$-fold map}, which is why the two band cases
        below need $z^{-2}$ and the two low-pass-like ones need only $z^{-1}$.
\end{itemize}}

\sbsr{0.56}{%
\lead{The four transformations \mk{(this is the table they want)}}
\par\smallskip
{\footnotesize
\begin{tabular}{@{}L{1.5cm}L{3.5cm}L{2.6cm}@{}}
\toprule
\hrow \thead{To} & \thead{Substitute $z^{-1}\to$} & \thead{Parameters} \\
Low pass & $\dfrac{z^{-1}-\alpha}{1-\alpha z^{-1}}$
  & $\alpha=\dfrac{\sin\frac{\omega_p-\omega'_p}{2}}
                  {\sin\frac{\omega_p+\omega'_p}{2}}$ \\[6pt]
High pass & $-\dfrac{z^{-1}+\alpha}{1+\alpha z^{-1}}$
  & $\alpha=-\dfrac{\cos\frac{\omega_p+\omega'_p}{2}}
                   {\cos\frac{\omega_p-\omega'_p}{2}}$ \\[6pt]
Band pass & $-\dfrac{z^{-2}-a_1z^{-1}+a_2}{a_2z^{-2}-a_1z^{-1}+1}$
  & $a_1=\dfrac{2\alpha k}{k+1}$, $a_2=\dfrac{k-1}{k+1}$ \\[6pt]
Band stop & $\dfrac{z^{-2}-a_1z^{-1}+a_2}{a_2z^{-2}-a_1z^{-1}+1}$
  & $a_1=\dfrac{2\alpha}{k+1}$, $a_2=\dfrac{1-k}{1+k}$ \\
\bottomrule
\end{tabular}}
\par\smallskip
{\color{sub}\footnotesize For the two band cases
$\alpha=\dfrac{\cos\frac{\omega_u+\omega_L}{2}}{\cos\frac{\omega_u-\omega_L}{2}}$, and
$k=\cot\frac{\omega_u-\omega_L}{2}\tan\frac{\omega_p}{2}$ for band pass,
$\tan\frac{\omega_u-\omega_L}{2}\tan\frac{\omega_p}{2}$ for band stop. $\omega_p$ is
the prototype's edge, $\omega'_p$ the new one.}}{%
\lead{Reading the table without mistakes}
\begin{itemize}
  \item \mk{Only the low-pass row keeps the sign}; the high-pass row carries a leading
        minus, and that minus is the whole reason the passband moves to $\omega=\pi$.
  \item The substitution goes into $H_{lp}(p)$ \textbf{written in $p^{-1}$}, not in $p$.
        Every published prototype already is.
  \item \textbf{Every zero at $z=-1$ becomes a zero at $z=+1$}, so a high pass built this
        way has $H(1)=0$ exactly: it kills DC. \mk{Free check on the answer.}
  \item $|\alpha|<1$ always. If it comes out bigger, the two edge frequencies have been
        swapped in the cosine.
  \item Both band rows raise the order: an $N$th-order low pass becomes a
        \textbf{$2N$th-order} band pass or band stop.
\end{itemize}}

\hr

\T{6.7 \quad Last-minute recall}

\sbs{%
\lead{Say these without thinking}
\begin{itemize}
  \item $\alpha=-20\log_{10}|H|$, and $\delta_p$ means an edge gain of $1-\delta_p$.
  \item $N=\dfrac{\log\left[(10^{0.1\alpha_s}-1)/(10^{0.1\alpha_p}-1)\right]}
                  {2\log(\Omega_s/\Omega_p)}$, \mk{round up}.
  \item $\Omega_c=\Omega_p\big/(10^{0.1\alpha_p}-1)^{1/2N}$.
  \item Bilinear: pre-warp $\Omega=\frac2T\tan\frac\omega2$, map back
        $s=\frac2T\frac{1-z^{-1}}{1+z^{-1}}$, and \mk{$T$ cancels}.
  \item Impulse invariance: $\Omega=\omega/T$, and
        $\frac{1}{s-p_k}\to\frac{1}{1-e^{p_kT}z^{-1}}$.
  \item Butterworth poles: radius $\Omega_c$, angles
        $\pi\frac{N+2k+1}{2N}$, spacing $\pi/N$.
  \item Chebyshev: $\varepsilon=\sqrt{10^{0.1\alpha_p}-1}$, $\Omega_c=\Omega_p$, poles on
        an ellipse.
  \item LP$\to$HP: $z^{-1}\to-\frac{z^{-1}+\alpha}{1+\alpha z^{-1}}$,
        $\alpha=-\frac{\cos\frac{\omega_p+\omega'_p}{2}}
                      {\cos\frac{\omega_p-\omega'_p}{2}}$.
  \item $\omega=2\pi f/f_s$ when the spec is in Hz.
\end{itemize}}{%
\lead{Mistakes that cost marks}
\begin{itemize}
  \item \mk{Forgetting to pre-warp} in a bilinear design, or pre-warping in an
        impulse-invariance one.
  \item Rounding $N$ down, or to nearest.
  \item Taking $\alpha_s$ as negative because the paper prints $-15$ dB. The formula
        wants $+15$.
  \item Reading $\delta_p=0.11$ as a gain of $0.11$. It is a \textbf{deviation}: the
        gain is $0.89$.
  \item Using all $2N$ roots of the Butterworth equation. \mk{Only the left half}.
  \item Substituting $z=e^{sT}$ into $H_a(s)$ for impulse invariance. That is not the
        method; only the \textbf{poles} move that way.
  \item Giving an even-order Chebyshev filter a DC gain of 1. It is
        $1/\sqrt{1+\varepsilon^2}$.
  \item Claiming impulse invariance can make a high-pass filter. \mk{It cannot.}
  \item Carrying $T$ through a bilinear design and getting a $T$-dependent answer.
        \mk{If $T$ survives, there is an arithmetic error.}
\end{itemize}}
